% Generated draft: requires mathematical and source review.
\documentclass[11pt]{article}
\usepackage[margin=25mm]{geometry}
\usepackage{fontspec}
\setmainfont{Latin Modern Roman}
\usepackage{amsmath,amssymb,mathtools,unicode-math}
\setmathfont{Latin Modern Math}
\usepackage{enumitem,needspace,xurl,hyperref}
\setlength{\parindent}{0pt}
\setlength{\parskip}{5pt}
\setlength{\emergencystretch}{3em}
\newcommand{\N}{\mathbb{N}}
\newcommand{\Z}{\mathbb{Z}}
\newcommand{\Q}{\mathbb{Q}}
\newcommand{\R}{\mathbb{R}}
\newcommand{\C}{\mathbb{C}}
\newcommand{\F}{\mathbb{F}}
\newcommand{\A}{\mathbb{A}}
\newcommand{\PP}{\mathbb{P}}
\newcommand{\E}{\mathbb{E}}
\newcommand{\eps}{\varepsilon}
\newcommand{\dd}{\,\mathrm{d}}
\begin{document}
\section*{UnsolvedMath ID: 30001968}
\textbf{Source problem code:} catalogue entry
\subsection{Definitions and mathematical statement}
\textbf{Modeling and Solving Bicycle Motion.} The supplied record presents the following methodological question: In developing a mathematical theory of bicycle motion, how should one choose coordinates, parameters, and variables; derive the differential equations; choose useful approximations; solve the equations analytically, numerically, or graphically; and draw practical conclusions?

\textbf{Unresolved formulation.} No specific mechanical model, configuration space, parameter domains, forces, constraints, differential equations, initial or boundary conditions, approximation criterion, or target conclusion is supplied. Consequently, the record does not determine a mathematical proposition with specified hypotheses and quantifiers. The tasks above are retained as a methodological inquiry; a precise theorem or conjecture cannot be extracted without additional source information.

The supplied statement-verification note characterizes the workshop discussion as historical, pedagogical, or methodological rather than an unresolved mathematical proposition. The record marks the statement as \texttt{unrecoverable}, provides no separate clean statement, and retains the original prompt for provenance. Its supplied \texttt{open} label is a curation classification, not independent verification of an open mathematical claim.
\subsection{Short English statement}
The source asks how to formulate, approximate, solve, and interpret a mathematical model of bicycle motion. It supplies neither a particular model nor a definite proposition, so a precise mathematical statement remains unavailable without further specifications. This is a statement-only draft requiring human mathematical review.
\subsection{Sources}
\begin{itemize}
\item Local UnsolvedMath registry, \texttt{data/unsolvedmath/problems.json}, permanent ID \texttt{30001968}, source problem code \texttt{catalogue entry}, title \emph{Modeling and Solving Bicycle Motion}; Oberwolfach Reports Open Problems collection. The supplied record assigns the year 2012 and names no individual proposer.
\item Supplied report citation: \emph{Explicit Versus Tacit Knowledge in Mathematics} (2012), Oberwolfach Reports, pp.~131--210, DOI 10.4171/owr/2012/04. Supplied source URL and literature reference labeled ``original Oberwolfach report'': \url{https://doi.org/10.4171/owr/2012/04}. No specific page within that range is supplied for this prompt.
\item Literature assessment dated exactly \texttt{2026-08-21}: the supplied curation labels the record \texttt{open} and describes the extracted text as methodological questions about modeling and solving mechanics problems, without a finite mathematical claim or established resolution. This is supplied dated curation, not an independently verified assessment. The supplied background expressly requests independent verification of time-sensitive status and cited claims.
\item The supplied statement-verification note characterizes the official history-of-mathematics workshop report as presenting a historical, pedagogical, or methodological discussion prompt. That characterization and the report attribution have not been independently checked in preparing this draft.
\end{itemize}
\end{document}
